
\documentclass{article}
\usepackage{colortbl}
\usepackage{makecell}
\usepackage{multirow}
\usepackage{supertabular}

\begin{document}

\newcounter{utterance}

\centering \large Interaction Transcript for game `referencegame', experiment `random\_grids', episode 3 with NEW\_Reflect\_Agent.
\vspace{24pt}

{ \footnotesize  \setcounter{utterance}{1}
\setlength{\tabcolsep}{0pt}
\begin{supertabular}{c@{$\;$}|p{.15\linewidth}@{}p{.15\linewidth}p{.15\linewidth}p{.15\linewidth}p{.15\linewidth}p{.15\linewidth}}
    \# & $\;$A & \multicolumn{4}{c}{Game Master} & $\;\:$B\\
    \hline

    \theutterance \stepcounter{utterance}  
    & & \multicolumn{4}{p{0.6\linewidth}}{
        \cellcolor[rgb]{0.9,0.9,0.9}{
            \makecell[{{p{\linewidth}}}]{
                \texttt{\tiny{[P1$\langle$GM]}}
                \texttt{You are given three grids, where each of them is 5 by 5 in size.} \\
\texttt{Grids have empty cells marked with "▢" and filled cells marked with "X".} \\
\texttt{Your task is to generate a referring expression that best describes the target grid while distinguishing it from the two other distractor grids.} \\
\texttt{The first grid is the target grid, and the following two grids are the distractors.} \\
\\ 
\texttt{Target grid:} \\
\\ 
\texttt{X X □ X □} \\
\texttt{□ X X □ X} \\
\texttt{X □ □ X □} \\
\texttt{□ X □ □ X} \\
\texttt{X □ X □ X} \\
\\ 
\texttt{Distractor grid 1:} \\
\\ 
\texttt{X X □ □ X} \\
\texttt{□ X X □ X} \\
\texttt{X □ X X □} \\
\texttt{□ X □ X X} \\
\texttt{X □ □ X X} \\
\\ 
\texttt{Distractor grid 2:} \\
\\ 
\texttt{X □ X X □} \\
\texttt{□ X □ X X} \\
\texttt{X X □ □ X} \\
\texttt{□ X X □ X} \\
\texttt{X □ X □ □} \\
\\ 
\texttt{Instruction: Describe the target grid.} \\
\texttt{Generate the referring expression starting with the tag "Expression: " for the given target grid. Omit any other text.} \\
            }
        }
    }
    & \\ \\

    \theutterance \stepcounter{utterance}  
    & \multicolumn{4}{p{0.45\linewidth}}{
        \cellcolor[rgb]{0.95,0.95,0.95}{
            \makecell[{{p{\linewidth}}}]{
                \texttt{\tiny{[P1$|$internal]}}
                \texttt{Reflections from previous episodes:} \\
\texttt{{-}} \\
\texttt{ Expression: The grid has exactly one filled cell in the first row, third column, and the pattern includes a diagonal of alternating filled and empty cells from top{-}left to bottom{-}right.} \\
\\ 
\texttt{Reflection: Failed. The expression was not precise enough to uniquely distinguish the target from both distractors. I focused on specific cell positions and a diagonal pattern, but distractors shared similar local features. Next time, I would systematically compare all three grids row by row, identify a minimal set of distinguishing cells (like a unique empty/filled cell present only in the target), and describe that exact combination concisely.} \\
\texttt{{-}} \\
\texttt{ I succeeded. The target grid had a distinctive pattern: a filled X in the center cell (row 3, column 3) and no fully alternating checkerboard like distractor 1, nor the diagonal clusters like distractor 2. My expression highlighted the central X and the asymmetry, which set it apart. Next time, I would similarly focus on unique structural features—like specific filled cells or pattern breaks—that clearly differentiate the target.} \\
\texttt{{-}} \\
\texttt{ Failed. I generated "Expression: The grid has a filled plus sign shape in the center." This was incorrect because it didn't uniquely distinguish the target. The target's key feature was the specific pattern of the **second row being "X □ □ X X"**, which neither distractor had. Next time, I will identify a single, unique row or column pattern that only the target possesses and describe that precisely.} \\
            }
        }
    }
    & & \\ \\

    \theutterance \stepcounter{utterance}  
    & \multicolumn{4}{p{0.45\linewidth}}{
        \cellcolor[rgb]{0.95,0.95,0.95}{
            \makecell[{{p{\linewidth}}}]{
                \texttt{\tiny{[P1$|$internal]}}
                \texttt{{[}attempt 1 – valid{]}} \\
\texttt{Expression: The grid has the second row as "X □ □ X X", which is unique to the target and not present in either distractor.} \\
            }
        }
    }
    & & \\ \\

    \theutterance \stepcounter{utterance}  
    & \multicolumn{4}{p{0.6\linewidth}}{
        \cellcolor[rgb]{0.9,0.9,0.9}{
            \makecell[{{p{\linewidth}}}]{
                \texttt{\tiny{[P1$\rangle$GM]}}
                \texttt{Expression: The grid has the second row as "X □ □ X X", which is unique to the target and not present in either distractor.} \\
            }
        }
    }
    & & \\ \\

    \theutterance \stepcounter{utterance}  
    & & & \multicolumn{2}{p{0.3\linewidth}}{
        \cellcolor[rgb]{0.9,0.9,0.9}{
            \makecell[{{p{\linewidth}}}]{
                \texttt{\tiny{[GM$|$GM]}}
                \texttt{Expression: The grid has the second row as "X □ □ X X", which is unique to the target and not present in either distractor.} \\
            }
        }
    }
    & & \\ \\

    \theutterance \stepcounter{utterance}  
    & & \multicolumn{4}{p{0.6\linewidth}}{
        \cellcolor[rgb]{0.9,0.9,0.9}{
            \makecell[{{p{\linewidth}}}]{
                \texttt{\tiny{[GM$\rangle$P2]}}
                \texttt{You are given three grids, where each of them is 5 by 5 in size.} \\
\texttt{Grids have empty cells marked with "▢" and filled cells marked with "X".} \\
\texttt{You are also given a referring expression that describes one of the given grids.} \\
\texttt{Your task is to select the grid that matches the given referring expression.} \\
\texttt{Generate only the number (in text) of the grid that matches the given expression by selecting first, second, or third.} \\
\\ 
\texttt{First:} \\
\\ 
\texttt{X X □ X □} \\
\texttt{□ X X □ X} \\
\texttt{X □ □ X □} \\
\texttt{□ X □ □ X} \\
\texttt{X □ X □ X} \\
\\ 
\texttt{Second:} \\
\\ 
\texttt{X X □ □ X} \\
\texttt{□ X X □ X} \\
\texttt{X □ X X □} \\
\texttt{□ X □ X X} \\
\texttt{X □ □ X X} \\
\\ 
\texttt{Third:} \\
\\ 
\texttt{X □ X X □} \\
\texttt{□ X □ X X} \\
\texttt{X X □ □ X} \\
\texttt{□ X X □ X} \\
\texttt{X □ X □ □} \\
\\ 
\texttt{Expression: The grid has the second row as "X □ □ X X", which is unique to the target and not present in either distractor.} \\
\texttt{Question: Which grid does the expression refer to?} \\
\texttt{Start with the tag "Answer: ", followed by your selection. Omit any other text.} \\
            }
        }
    }
    & \\ \\

    \theutterance \stepcounter{utterance}  
    & & & \multicolumn{4}{p{0.45\linewidth}}{
        \cellcolor[rgb]{0.95,0.95,0.95}{
            \makecell[{{p{\linewidth}}}]{
                \texttt{\tiny{[P2$|$internal]}}
                \texttt{Reflections from previous episodes:} \\
\texttt{{-}} \\
\texttt{ Failed. I misinterpreted "vertical stripe of X's in the middle three columns" as all cells in those columns being X's. Grid 1 has the correct empty left/right columns, but the middle columns are not completely filled with X's. I should have checked each column for consistency with the description instead of making a quick visual assumption. Next time, verify each column individually against the stated pattern.} \\
\texttt{{-}} \\
\texttt{ Failed. I misinterpreted "first row starting with an empty cell" as describing the whole grid's first row, but the expression was describing the *target grid's* first row. The correct grid (First) actually starts with an empty cell (▢). I focused too much on the diagonal description and missed this key initial condition. Next time, I will check each condition in the expression sequentially against all grids before deciding.} \\
\texttt{{-}} \\
\texttt{ Answer: Third} \\
\\ 
\texttt{Reflection: Succeeded. I carefully parsed the expression, checking each grid for the specific filled cell at (1,3) and the described diagonal pattern from (2,1) to (4,3) with alternating cells and an empty center. Only the third grid matched all conditions. Next time, I will continue to methodically verify each stated feature against the grids to ensure accuracy.} \\
            }
        }
    }
     \\ \\

    \theutterance \stepcounter{utterance}  
    & & & \multicolumn{4}{p{0.45\linewidth}}{
        \cellcolor[rgb]{0.95,0.95,0.95}{
            \makecell[{{p{\linewidth}}}]{
                \texttt{\tiny{[P2$|$internal]}}
                \texttt{{[}attempt 1 – valid{]}} \\
\texttt{Answer: third} \\
            }
        }
    }
     \\ \\

    \theutterance \stepcounter{utterance}  
    & & & \multicolumn{4}{p{0.6\linewidth}}{
        \cellcolor[rgb]{0.9,0.9,0.9}{
            \makecell[{{p{\linewidth}}}]{
                \texttt{\tiny{[GM$\langle$P2]}}
                \texttt{Answer: third} \\
            }
        }
    }
     \\ \\

    \theutterance \stepcounter{utterance}  
    & & & \multicolumn{2}{p{0.3\linewidth}}{
        \cellcolor[rgb]{0.9,0.9,0.9}{
            \makecell[{{p{\linewidth}}}]{
                \texttt{\tiny{[GM$|$GM]}}
                \texttt{third} \\
            }
        }
    }
    & & \\ \\

\end{supertabular}
}

\end{document}
